\section{Introdução}

Detectar resíduos em ambientes externos é diferente de reconhecer itens em condições controladas: os objetos podem ser pequenos, deformáveis, transparentes, ocluídos ou visualmente confundidos com o fundo. O TACO (\textit{Trash Annotations in Context}) foi criado para representar esse cenário, reunindo fotografias de ambientes diversos e uma taxonomia hierárquica de resíduos \cite{proenca2020taco}. Entretanto, sua versão oficial contém apenas 1.500 imagens e 4.784 instâncias distribuídas em 60 categorias, combinação que impõe simultaneamente poucos dados e forte desbalanceamento.

Em detecção multiclasse, uma caixa espacialmente correta é considerada erro quando recebe a classe errada. Assim, o desempenho pode diminuir tanto por falhas de localização quanto pela exigência de distinguir categorias visualmente próximas e pouco representadas. A granularidade da taxonomia determina quanta informação semântica se exige do modelo, mas comparações entre estudos costumam misturar arquiteturas, partições e agrupamentos diferentes. Isso impede atribuir a variação de desempenho apenas aos rótulos.

Este trabalho investiga: \emph{como a granularidade das classes afeta a detecção de resíduos em um conjunto pequeno e de cauda longa, e quanto da perda decorre de localização versus discriminação semântica?} Para isolar essa variável, fixamos imagens, caixas, partição e protocolo de treino, alterando somente a taxonomia: Fine (60 classes oficiais), Material (seis macroclasses) e Binary (uma classe \textit{Litter}). Empregamos o YOLO26n, variante compacta da família YOLO26 \cite{jocher2026yolo26}, sem modificar sua arquitetura.

As contribuições são: (i) comparação controlada dos três níveis com três sementes; (ii) avaliação hierárquica das \emph{mesmas} predições Fine, remapeadas para material e para classe única; (iii) intervalos de confiança por bootstrap pareado por imagem, repetidos independentemente para cada semente; e (iv) análises multissemente de cauda longa e uma decomposição auditável dos erros. Todo número apresentado é produzido por scripts a partir das anotações, configurações, checkpoints ou predições salvas.
